% Discussion (v2, trade-off framing): one cause behind the gains and the losses, what the control
% changes, and where the design should pay. Caveats are stated once, in Limitations.
\section{Discussion}
\label{sec:discussion}

\paragraph{One account of both sides of the trade-off.}
A claim replaces its passage as the generator's input. Because claims accumulate, a page
can grow with the story and keep what it said, which a page rewritten under a fixed budget
cannot (Table~\ref{tab:gains}). Because a claim carries a short quote rather than its
surrounding discourse, the filter tests the one property the claim retains exactly, the
volume of its evidence; speaker, mood, which
namesake a mention denotes and whether two characters share an episode are held weakly or
not at all. Hence a prediction became a marriage, entities took their shape from later
volumes, and true lines became unsupported where published (Tables~\ref{tab:kinds}
and~\ref{tab:identity}). B4, reading the same graph through retrieved facts, kept a hedge
that deterministic rendering dropped but lost the stability rendering gave: both effects
sit at publication, not only at extraction, and stage ablations can separate them.

\paragraph{Reading must be shown, not assumed.}
The renamed-text control removes the floor that a known work gives a generator reading
nothing, and changes three conclusions (Section~\ref{sec:results-decon}). Evaluations on
published narratives should report a recognition check and an aligned
treated condition, and audit the treatment (Appendix~\ref{app:contamination}).

\paragraph{Where the trade-off should pay.}
Accumulation matters most where rewriting each page is costly or impossible: series longer
than one context window, and characters whose record outgrows a page budget. The failures
are repairable without giving up accumulation: carry speaker and modality in the claim
schema, validate entity groupings against the prefix, and check each final assertion in
its page context against its citations (Appendix~\ref{app:casebook}).
